\section{Formal Grammar Specification}\label{sec:formal-grammar}

The proposed DSL is defined by a formal grammar that ensures syntactic consistency and enables the generation of an Abstract Syntax Tree (AST).
This grammar is divided into a lexical specification (the scanner) and a context-free grammar (the parser).

\subsection{Lexical Grammar}

The lexical layer tokenizes the raw source text into a stream of symbols.
A notable feature of the lexical grammar is the inclusion of musical literals as first-class tokens.
For example, a note literal is defined by the regular expression \texttt{[A-G][\#b]?[0-8]}, allowing the scanner to directly recognize pitches such as \texttt{F\#4} or \texttt{Bb2}.
Keywords are reserved for musical structural elements (\texttt{track}, \texttt{pattern}, \texttt{voice}), temporal modifiers (\texttt{play}, \texttt{from}, \texttt{rest}), and standard programming constructs (\texttt{let}, \texttt{for}, \texttt{if}).

\subsection{Syntactic Grammar}

The syntactic grammar follows an EBNF (Extended Backus-Naur Form) structure.
It is designed to enforce the hierarchical relationship between musical components.
The root of a program consists of a header (defining metadata like tempo) followed by top-level declarations.

\begin{figure}[htbp]
    \centering
    \begin{small}
    \begin{verbatim}
<program>    ::= <header> <top-item>*
<top-item>   ::= <let-stmt> | <pattern-def> | <track-decl>
<track-decl> ::= "track" IDENT? ("using" <inst>)? "{" <track-body> "}"
<track-body> ::= (<stmt> | <pattern-def> | <voice-decl>)*
<pattern>    ::= "pattern" IDENT "(" <params> ")" "{" <stmt-list> "}"
<play-stmt>  ::= "play" <play-target> ";"
<play-target>::= <source> (":" <expr>)? ("from" <expr>)?
    \end{verbatim}
    \end{small}
    \caption{Simplified EBNF representation of the core DSL grammar.}
    \label{fig:simplified-grammar}
\end{figure}

The grammar allows for both braced and unbraced statement bodies in control flow structures, providing a balance between clarity and brevity.
Crucially, the grammar restricts certain constructs based on their context; for instance, while a \texttt{track} can contain multiple \texttt{voice} blocks, a \texttt{voice} cannot nest another \texttt{track}, enforcing a one-to-one mapping between tracks and instrument identities.
The expression sub-grammar supports a full range of arithmetic and logical operators, including a ternary conditional operator, enabling complex musical logic to be expressed inline.
